\section{An actual auxiliary ordinary seed on the certified host tower}
\label{sec:cubic-auxiliary-host-seed}

This section starts with the actual form of AC-G119, rather than an
original-level Type II point.  That distinction is essential: the
order-$269$ inertia in AC-G119 violates the original residual-kernel
killing condition before restriction.  We remove that inertia on an
explicit real cyclic extension before constructing the definite
unitary ordinary module.  No original $\mathrm{Dat.IsPt}$ is asserted.

Write $c_p$ for the actual cyclotomic character, so that
$c_p(\operatorname{Frob}^{\rm geom}_r)=r^{-1}$.
All Frobenius polynomials below use geometric Frobenius.
Let $\chi$ be the fixed cubic character of $G_K$, with
$K=\mathbb Q(\sqrt{-3})$, and put $\eta=\eta_K$.

\begin{theorem}[AC-G120-auxiliary-host: an actual ordinary irreducible host seed]
\label{thm:G120}
Set
\[
 \begin{gathered}
 p=269,\quad q_0=3229,\quad \psi=\chi_{-4},\quad
 e\equiv267\pmod{268},\quad e\ge1,\quad k=2e+1,\\
 K'=\mathbb Q(\sqrt{-23}),\quad F_1=KK',\quad
 F_1^+=\mathbb Q(\sqrt{69}),\quad E=\mathbb Q(\sqrt5),\\
 F_{\rm ab}=F_1^+E,\qquad H_{\rm ab}=F_1E.
 \end{gathered}                                                     \tag{G120.1}
\]
Let $T_\theta/\mathbb Q$ be the real cyclic degree-$p$ field of
AC-G119, and let $f_0,V$ be its actual ordinary form and normalized
rank-two representation.  In particular
\[
 \begin{gathered}
 \operatorname{level}(f_0)=48q_0^2,\quad
 \operatorname{neb}(f_0)=\eta,\quad \det V=\eta,\\
 \bar V^{\rm ss}=\eta\bar\psi\bar c_p^{\,e}
                         \oplus\bar\psi^{-1}\bar c_p^{-e},\qquad
 V|_{I_{q_0}}^{\rm ss}=\Theta^{1/2}\oplus\Theta^{-1/2},\quad N=0.
 \end{gathered}                                                     \tag{G120.2}
\]
The finite character $\Theta$ cuts out $T_\theta$.
Put
\[
                   F=F_1T_\theta,\qquad F^+=F_1^+T_\theta.
                                                                    \tag{G120.3}
\]
Retain the same published ordinary-control, Galois-attachment, cyclic
base-change, endoscopic-transfer, local--global compatibility and local
type inputs explicitly used in AC-G75 and AC-G78--AC-G83.  Their
arithmetic starting-form hypotheses are supplied here by AC-G119 and
the proof below; no original-level endpoint-existence premise remains.

For each $j\in\{1,2\}$ there is a rational prime $q\ne q_0$, split
completely in $F$, with
\[
 q\equiv-1\pmod p,\quad v_p(q+1)=1,\quad
 \chi(\operatorname{Frob}_q)=1,\quad q\text{ inert in }E,
                                                                    \tag{G120.4}
\]
and there is a totally real dihedral field $D_q/\mathbb Q$ whose
quadratic subfield is $E$, with $[D_q:E]=p$, ramified only above $q$
over $E$ and with inertia order $p$ there.
There is an actual ordinary automorphic representation on the definite
unitary group for
\[
 H=FD_q,
 \qquad H^+=F_1^+T_\theta D_q,
 \qquad [H:\mathbb Q]=8p^2,
 \qquad [H^+:\mathbb Q]=4p^2.                       \tag{G120.5}
\]
Its weight is the original parallel symmetric weight
$\mu(e)=(e-1,0,1-e)$.  Its normalized Galois representation $r_H$
has the following properties:
\begin{gather}
 r_H\text{ is absolutely irreducible},\notag\\
 \bar r_H^{\rm ss}
  =\left(\chi^j\bar\psi\bar c_p^{\,e}\right)
       \oplus1\oplus
    \left(\chi^j\bar\psi^{-1}\bar c_p^{-e}\right)\bigm|_{G_H};
                                                               \tag{G120.6}\\
 N_v=0,\qquad
 r_H|_{I_v}^{\rm ss}=(\psi,1,\psi^{-1})|_{I_v}
                      \quad(v\nmid p).                         \tag{G120.7}
\end{gather}
Thus $r_H$ is unramified above both $q_0$ and $q$.
At every place over $p$, it is potentially crystalline, has $N=0$,
and has an ordered invariant flag with graded characters
\[
 u_1\chi^j c_p^e,\qquad u_2,\qquad u_3\chi^j c_p^{-e},
               \qquad u_i\text{ unramified}.                    \tag{G120.8}
\]
The completion $H_w/\mathbb Q_p$ is unramified of degree $538$.
This is a statement about the local number field, not a change of the
original coefficient DVR.  The construction permits finite value-field
extensions, as the cited ordinary and local-type constructions do.
The extension $H/H^+$ is unramified at every finite place, all places
over $p$ split relatively, and no completion above $p$ contains $\mu_p$.

The same rank-two input restricted to $H$ is absolutely irreducible
and has finite away-$p$ inertia $(\psi,\psi^{-1})$, with $N=0$.
This is an actual host-level killing assertion, not an original-level
$\mathrm{Dat.Kill}$ assertion.
If $\Omega/\mathbb Q$ is a finite Galois extension of degree prime
to $p$, contains $K$, and is unramified at $5$ and $23$, then
\[
                         H\cap\Omega=K.                         \tag{G120.9a}
\]
In particular this holds for the full original finite residual data.

The independent arithmetic certificate of AC-G103 on $H_{\rm ab}$
and the rank-transport theorem AC-G100 apply to
\[
               H_{\rm ab}\subset H_{\rm ab}T_\theta\subset H.
                                                                    \tag{G120.9}
\]
Consequently, for the actual adjacent coefficient
$\alpha=\chi^{-j}\bar\psi^{-1}\bar c_p^{-e}$,
\[
               r_{\alpha,H}\le\tfrac34[H^+:\mathbb Q].             \tag{G120.10}
\]
This compatibility uses the unconditional finite certificate only;
it does not reuse the empty original-level seed premise in the
automorphic clause of AC-G103.
The selected-place extensions of AC-G92 and the cyclotomic steps of
AC-G93-fill can subsequently be applied with their stated hypotheses.
The theorem itself does not assert the native support map, a uniform
even-prime support ideal, any prescribed Hecke-curve containment,
or an original-level deformation point.
\end{theorem}

\begin{proof}
\emph{1. The actual rank-two input and its restrictions.}
AC-G119 supplies (G120.2) from an integral Eisenstein cusp lift and
Deligne--Serre eigenvalue lifting, including its exact primitive level.
In particular the whole fixed-level ordinary Hecke branch is
unramified at $5$ and $23$.
The representation $V$ is the unitary cyclotomic normalization of
the actual primitive $f_0$, not a coefficientwise nonprimitive twist.
At $p$ the untwisted cohomological representation has two different
inertial characters $1,c_p^{-2e}$ and one uniquely determined
unramified ordinary subquotient.
If $V|_{G_K}$ were reducible, Clifford theory for the quadratic
extension would give $V\simeq V\otimes\eta$.
This would preserve that unique unramified ordinary subquotient.
But $p$ is inert in $K$, so its unramified quadratic twist has
Frobenius value multiplied by $-1$, which cannot preserve a nonzero
characteristic-zero value.  Thus $V|_{G_K}$ is absolutely irreducible.
The same argument may be made on the cohomological representation
before its common scalar normalization.

The further quadratic restriction to $F_1$ remains irreducible.
Otherwise it gives the self-twist by $K'/\mathbb Q$ on $G_K$;
that quadratic character is ramified at $23$, while $V$ is unramified
there.  An unramified representation cannot equal its ramified
scalar quadratic twist.  Restriction to $F=F_1T_\theta$ also remains
irreducible.  Indeed a cyclic prime-$p$ quotient with an irreducible
rank-two representation can cause reducibility only through a
nontrivial quotient-character self-twist; determinants would then
give $\nu^2=1$, impossible for a nontrivial character of order $p$.
Equivalently, the translates of a simple restricted module give a
semisimple restriction, and Frobenius reciprocity computes its
commutant as the sum of these self-twist spaces.

Choose a minimal ordinary cuspidal Hecke branch through $f_0$.
As in AC-G78, after fixing its Teichmuller component it is a complete
domain $R$ finite torsionfree over $\mathcal O[[W]]$, with the
normalized continuous rank-two determinant and trace $\mathcal T$.
The tame nebentype is exactly $\eta$ on this domain, even though
the tame diamond group has order divisible by $p$.  If a diamond
value $a\in R$ has finite order $d$ and specializes to
$\zeta=\eta(\text{diamond})$, factor
$a^d-1=(a-\zeta)Q(a)$.  Its specialization has
$Q(a)(f_0)=d\zeta^{d-1}\ne0$ in characteristic zero.
Since $R$ is a domain, $a=\zeta$.  This is the AC-G88
domain argument; no tame diamond idempotent of order $p$ is used.
It has determinant $\eta$, and the parameter is chosen so that the
classical point of weight $2w+1$ has $1+W=(1+p)^w$.
The arithmetic points with $w>0$ and $w\equiv e\pmod{268}$ are dense.
The no-$\eta$-self-twist condition holds on a dense open subset:
some $\mathcal T(a)-\eta(a)\mathcal T(a)$ is nonzero at the starting
point.  The fresh-$23$ argument works for every classical point of
this fixed tame level.  The subsequent cyclic-$p$ restrictions
preserve rank-two irreducibility by the determinant argument.
Thus the required endpoint base changes are cuspidal on a dense
set of actual classical points.

\emph{2. Remove the auxiliary type on the whole branch first.}
Apply Saha, \emph{Purity for families of Galois representations},
Section 6.1.1, Theorem 6.2, to this minimal branch and $q_0$.
The complete statement compares the generic Frobenius-semisimple
Weil--Deligne representation with every arithmetic specialization.
The special representation alternative has nonzero monodromy at
every arithmetic point, and is excluded by (G120.2).
In the remaining alternatives the finite inertia representation
is constant.  To see the precise meaning here, its matrices factor
through a fixed finite group.  Its character values are algebraic
constants, sums of roots of unity.  A characteristic-zero
specialization is injective on these constants; after a finite
coefficient extension on a component through $f_0$, its inertia
character is exactly $\Theta^{1/2}+\Theta^{-1/2}$.
Finite-group semisimplicity identifies the representation, and
monodromy is zero.  This reasoning does not divide by $p$ in an
integral automorphic module and uses no residual irreducibility.

Since $\Theta|_{G_{T_\theta}}=1$, every such classical endpoint
restricted to $F$ is unramified at all places above $q_0$.
This is established before choosing the definite unitary level.
We therefore choose that level hyperspecial at every $q_0$-place;
there is no order-$p$ character projector to be applied afterward.
At $2,3$ Saha gives the same fixed types as at the initial form.
After restriction to $K$, the $3$-type $1\oplus\eta$ is trivial,
and the $2$-type is $\psi\oplus\psi^{-1}$.
The analytic weight normalization is unramified at these primes.
The cubic character is trivial locally at $3$: $269\equiv-1\pmod9$
is a cube in $\mathbb Q_3$, as in AC-G80.
Thus the full rank-three endoscopic parameter has the finite
away-$p$ type $(\psi,1,\psi^{-1})$ and $N=0$ at every fixed bad place.

\emph{3. The enlarged CM field and its local degrees.}
The discriminants of $K,K'$ are $-3,-23$ and are coprime.
The biquadratic discriminant formula gives
$|\operatorname{disc}F_1|=69^2$ and
$\operatorname{disc}F_1^+=69$, so $F_1/F_1^+$ is unramified
at every finite place.  We have
\[
 -23\equiv1\pmod8,\quad -23\equiv1\pmod3,\quad
 28^2\equiv-23\pmod{269},\quad q_0\equiv9\pmod{23}.
                                                                    \tag{G120.11}
\]
Hence $2,3,p,q_0$ split in $K'$; for $q_0$ one also uses
$(-23/q_0)=(q_0/23)$ and $q_0\equiv1\pmod3$.
All ramified places of the endpoint and of $\chi$ are relatively
split for $F_1/F_1^+$, and this remains true after real base change.
The field $T_\theta$ is real abelian of degree $p$, ramified only
at $q_0$.  It is disjoint from $F_1$ by their coprime degrees.
Therefore $F/F^+$ is finite-unramified, $F=F^+K'$, and all
coefficient-prime places split relatively.

AC-G119's exact discrete-log calculation gives
$\theta(p)=\zeta_p^{175}$, of order $p$.
Thus $(T_\theta)_p/\mathbb Q_p$ is unramified of degree $p$.
The completion of $F_1$ above $p$ is the unramified quadratic
field $K_p$, so
\[
                    [F_w:\mathbb Q_p]=2p=538,\qquad e(F_w/\mathbb Q_p)=1.
                                                                    \tag{G120.12}
\]
The finite cubic inertia of $\chi$ is unchanged.
Neither this local field nor any finite extension of $\mathbb Q_p$
of degree a power of $2$ times a power of $269$ can contain
$\mu_{269}$: its degree would
then be divisible by $268=4\cdot67$.
The coefficient ring used to realize embeddings and finite types
may be enlarged separately; (G120.12) is not a coefficient-ring
ramification statement.

\emph{4. Actual endoscopic ordinary points on this field.}
For each dense classical point use
\[
          \pi_w=(\operatorname{BC}_F(f_w)^{\rm unit}\otimes\chi^j)
                         \boxplus1 .                              \tag{G120.13}
\]
It is regular algebraic and conjugate self-dual with weight
$\mu(w)=(w-1,0,1-w)$; its endpoint is cuspidal by Step 1.
Take $F_0=K'$ in Thorne's setup.  The shifted pair of parallel
archimedean weights in Theorem 4.5 occurs $2p$ times above
each embedding of $F_0$, an even number.  All ramified places
are relatively split by Steps 2--3.  Theorem 4.5 therefore
gives the actual definite-unitary packet with the exact split
local components.  Proposition 6.5 gives its actual ordinary
refinement.  Its proof and hypotheses impose no bound on the
unramified local degree in (G120.12).

Use the integral projectors of AC-G80 for the finite away types.
Their images have order prime to $p$, since the order-$p$ type
at $q_0$ was already killed.  Repeated characters are selected
by the full maximal compact in their Levi block, hence impose
$N=0$, not only an Iwahori condition.  Include the auxiliary
neat-level restoring projector from the beginning, exactly as
in AC-G72/G80.  Such a place may be chosen split in $F$ with
residue modulo $p$ of order greater than $3$; then its finite
$\mathrm{GL}_3$ quotient has order prime to $p$.  The intersection
$F\cap\mathbb Q(\mu_p)=\mathbb Q$ follows since $F$ is unramified
at $p$ and every nontrivial subfield of $\mathbb Q(\mu_p)$ is
ramified there.  Chebotarev thus supplies this place outside
any fixed finite set.  This is no post hoc averaging argument.

At $p$ retain the full finite unit-torus character projector
of AC-G79.  For $L=F_w$, put $f=[L:\mathbb Q_p]=538$.
The T-equivariant invertible-$U$ to Jacquet map of
Clozel--Thorne I, Lemma 2.5, and lowest-weight evaluation give
the full unit character
\[
          \delta_{\rm full}=\alpha\,w_0(\mu)^{-1}.
                                                                    \tag{G120.14}
\]
The ordered inducing slopes are $(-fw,0,fw)$.
The ordinary scalar $U$-eigenvalue has valuation
$2f-2fw+2f(w-1)=0$; all other permutations have strictly
positive rearrangement excess.  This is exactly the AC-G79
calculation with the actual value of $f$, not a theorem
transported from $\mathbb Q_p$.
The fixed tame tuple is $(\chi^j,1,\chi^j)$.
For $w\equiv e\pmod{268}$ the parallel algebraic characters
have the same Teichmuller restriction, since it factors through
the residue norm to $\mathbb F_p^\times$.
Thus all the dense points in (G120.13) belong to one integral
finite-torus summand.

The finite-depth freeness, trace/norm compatibility and actual
continuous-weight specialization of AC-G78-module apply to
this field without change.  The product of the auxiliary,
away-type and finite-torus projectors commutes with these
maps, the ordinary projector and all away Hecke operators.
It therefore gives the same finite free ordinary module on
the full weight algebra.  The parallel branch character is
\[
 \lambda_R(t)=
 \frac{\nu_R\bigl(\prod_{v\mid p}N_{F_{\widetilde v}/\mathbb Q_p}
                         (t_{v,1}/t_{v,3})\bigr)}
      {\prod_{v\mid p}N_{F_{\widetilde v}/\mathbb Q_p}
                         (t_{v,1}/t_{v,3})},\qquad \nu_R(1+p)=1+W.
                                                                    \tag{G120.15}
\]
This is the actual parallel embedding character also for (G120.12).
The dense points give a map from the actual finite Hecke image
to $R$.  As in AC-G78, every operator relation evaluates to zero
on those points and hence in the domain $R$.
Upon arbitrary specialization the polynomial-law Cayley--Hamilton
identity kills the kernel of the specialized matrix action.
The resulting character is consequently on the actual specialized
ordinary module, not just a formal interpolation ring.

\emph{5. The Kummer image after adjoining $T_\theta$.}
Let $u$ be the square of a fundamental unit of $E$ and put
\[
                   M=E(\mu_p,u^{1/p}),\qquad
                   A=\operatorname{Gal}(M/E(\mu_p)).
                                                                    \tag{G120.16}
\]
This is a Galois extension of $\mathbb Q$, since the quadratic
conjugate of $u$ is $u^{-1}$, and $A$ is trivial or cyclic of
order $p$.  It is unramified at $q_0$: the unit Kummer extension
is unramified away from $p$, and $E$ ramifies only at $5$.
The field $F_1M$ is therefore unramified at $q_0$, whereas
$T_\theta/\mathbb Q$ is totally ramified there of degree $p$.
Hence
\[
               T_\theta\cap F_1M=\mathbb Q,
 \qquad \operatorname{Gal}(FM/F)
                  \simeq\operatorname{Gal}(F_1M/F_1).             \tag{G120.17}
\]
The latter image contains $A$, because its index in
$\operatorname{Gal}(M/\mathbb Q)$ divides $[F_1:\mathbb Q]=4$.
Thus $G_F$ still surjects onto $A$.  This proves the needed
replacement for the old degree-four image argument; the
degree $4p$ of $F$ alone would not prove it.

Choose $h\in I_5$ with $\eta_E(h)=-1$.
The fields $F_1,T_\theta$ are unramified at $5$, so $h\in G_F$.
The entire fixed-level branch is unramified at $5$, hence
\[
 \mathcal T(ha)=\mathcal T(a)\quad(a\in G_F),\qquad
                  c_p(h)=\chi(h)=\psi(h)=1.                       \tag{G120.18}
\]
Choose $t\in I_p\cap G_F$ with $c_p(t)=-(1+p)$ and $\chi(t)=1$.
The host is unramified at $p$, so its inertia is unchanged.
The tame orders $268$ and $3$ are coprime, and the wild
cyclotomic value is unrestricted, so this simultaneous choice
exists.  Also $\eta_E(t)=1$ and $\psi(t)=1$.
The odd-$e$ ordinary trace identity is
\[
              \mathcal T(t)=-((1+W)+(1+W)^{-1}).                  \tag{G120.19}
\]

Conjugation on $A$ has character $\eta_E\bar c_p$;
both $h$ and $t$ therefore act by inversion.
For $g=ht$ one has $g^2\in A$ and $g$ centralizes $A$.
By (G120.17), lift any $a_0\in A$ to $n\in G_F$.
Then $(nhn^{-1}t)^2=a_0^4g^2$ in $\operatorname{Gal}(M/\mathbb Q)$.
Since fourth powering is invertible on $A$, choose $a_0$ to
make this square one.  Cyclicity of trace gives
$\mathcal T(nhn^{-1}b)=\mathcal T(b)$ for every $b\in G_F$,
so this conjugation preserves the exact branch identity.
Put $g'=nhn^{-1}t$.  It has all the original analytic target
values and $(g')^2|_M=1$.

Use the complete exact-root argument of AC-G78, or its fixed
normalization refinement AC-G84, with this actual $g'$.
Add $F,M$, the residual characters and $c_p\bmod p^2$ to its
finite Chebotarev quotient.  The specified element proves
compatibility of all conditions in one quotient.
It gives a good rational $q$ satisfying (G120.4), with
$\psi(q)=1$, and with the norm-one unit $u$ a $p$th power
modulo the inert $E$-prime above $q$.
Indeed the Frobenius on $M$ squares to one, while its image
on $E(\mu_p)$ has order two, so that prime splits in the
Kummer extension.  The congruence
$c_p(\operatorname{Frob}_q)=q^{-1}\equiv-(1+p)\pmod{p^2}$
gives $v_p(q+1)=1$.
At an actual nearby $R$-point the rank-two roots are exactly
$q,q^{-1}$; the unitary rank-three roots are $q,1,q^{-1}$,
and Thorne's common $c_p^{-1}$ normalization gives $q^2,q,1$.
Step 4 places this point in the actual ordinary Hecke support.

\emph{6. Finite type, irreducibility and the ordered $p$-flag.}
The residual genericity witness of AC-G78 still exists.
The field $F/\mathbb Q$ is abelian, while the fixed cubic
splitting field has maximal abelian subfield $K$.
Thus its cubic value can be varied while fixing $F$ and all
rational characters.  Starting with $p$-inertia of cyclotomic
residue $-1$ yields normalized roots
$-\zeta_3,1,-\zeta_3$ at a good split witness prime.
Their ratios are never $-1$; this is the literal residual
genericity condition in Thorne's Theorem 5.2.

Apply the four actual elliptic tests of AC-G75, retaining all
three commuting projectors from Step 4, including the transformed
finite-torus summand in its fourth test.
They supply an ordinary Steinberg constituent at the original
classical weight.  At full $q$-level the lattice is projective
over $\mathcal O[\mathrm{GL}_3(\mathbb F_q)]$ by the same
free double-coset argument as AC-G72/G78-type.
Every constituent is generic by the retained residual witness
and Thorne's post-Corollary-4.3 criterion.
A spherical generic constituent has finite Steinberg and trivial
multiplicities both one; nonspherical constituents have trivial
multiplicity zero, and the actual smooth Steinberg contributes
one to the positive difference.
The proved Brauer identity
$[\bar\tau_{\vartheta}]=[\overline{\rm St}_3]-[1]$
therefore gives actual occurrence of the irreducible finite
type $\tau_{\vartheta}$, with $\vartheta$ regular of order $p$.
This uses projectivity together with the positive multiplicity
difference, not Steinberg presence alone.

AC-G72-local identifies the $q$-parameter as a tame irreducible
two-dimensional block with characters $\vartheta,\vartheta^q$
and one unramified line, with $N=0$.
Its residual roots are $-1,1,-1$.  The block's Frobenius roots
are $a,-a$, so the remaining line has residual value $-1$.
Among the three global residual characters in (G120.6), only
the middle one is individually conjugate self-dual:
$\bar\gamma^2$ is nontrivial on $p$-inertia, and the cubic
character cannot cancel its prime-to-three rational inertia.
A global semisimple $2+1$ decomposition would therefore have
its one-dimensional summand reducing to the middle character,
contrary to the local value $-1$.  Three global lines cannot
contain the irreducible local two-dimensional block.
Thus the normalized rank-three representation on $G_F$ is
absolutely irreducible.

Finally apply (G120.14) at the classical weight $e$.
The actual smooth finite inertial characters are
$(\chi^j,1,\chi^j)$ with slopes $(-fe,0,fe)$, $f=538$.
Monodromy commutes with inertia and connects slopes differing
by $f$.  The middle inertial character is different, and the
equal endpoint characters have slope gap $2ef\ne f$.
Hence $N=0$.  Local--global compatibility gives de Rhamness
and this full Weil--Deligne parameter.
After a finite semistabilizing extension, the parallel
filtration jumps and normalized Newton slopes are $(-e,0,e)$.
For each initial Newton-slope subspace, its Hodge number is
at least the sum of the corresponding smallest Hodge jumps,
equal to its Newton number; weak admissibility supplies the
opposite inequality.  Equality holds at each embedding.
These subspaces are weakly admissible, hence give invariant
subrepresentations by the Colmez--Fontaine equivalence.
Their unique slope descriptions make the flag descend.
The rank-one quotients have exactly (G120.8), since removal
of their parallel cyclotomic powers and finite inertia leaves
weight-zero unramified characters.  This repeats the full
AC-G79 argument for $f=538$, with no $\mathbb Q_p$-only input.

\emph{7. Kill the new $q$-type and identify the actual final host.}
The ray construction AC-G83 applies to $E,q,u$ from Step 5.
The class number of $E=\mathbb Q(\sqrt5)$ is one: its
Minkowski bound $\sqrt5/2<2$ leaves only the unit ideal.
Thus the ray $p$-part is exactly $C_p$ and gives $D_q/E$
of degree $p$, with all the assertions in (G120.4)--(G120.5).
The field $F/\mathbb Q$ is abelian.  Every abelian quotient
of the dihedral group of $D_q/\mathbb Q$ factors through
its quadratic quotient $E$.  But $E\not\subset F$, by the
fresh ramification at $5$.  Therefore $F\cap D_q=\mathbb Q$
and the displayed degree of $H$ follows.
The tower $F\subset FE\subset H$ has degrees $2,p$.
The field $D_q$ is totally real, so $H$ is CM and $H/H^+$
is the base change of $F_1/F_1^+$, hence finite-unramified.

At $q$, $E$ supplies the unramified quadratic extension,
and $D_q/E$ has tame inertia order $p$.  The two finite
characters of order $p$ in the $q$-parameter are thereby
killed.  The unramified line and $N=0$ persist.
There is no other ramification in $D_q/E$; at $5$ the
original parameter was unramified, so adjoining $E$ adds
no representation ramification.  Every old type simply
restricts.  This proves (G120.7), including at $q_0$.

At $p$, $E$ splits and $D_q/E$ is unramified of local
degree dividing $p$.  The field in (G120.12) already contains
the unique unramified degree-$p$ extension of $\mathbb Q_p$.
Thus the final local field $H_w$ is still the field in
(G120.12); (G120.8) and ordinarity persist exactly.
Rank-three irreducibility survives the cyclic steps $2,p$:
any nontrivial quotient-character self-twist would satisfy
$\nu^3=1$ by determinants, impossible for those orders.
The same argument with exponent two proves the rank-two
assertion at the cyclic-$p$ steps.  For its intervening
quadratic step $F\subset FE$, use instead the fresh prime
$5$: $V|_{G_F}$ is unramified there, whereas $FE/F$ is
ramified there.  The quadratic self-twist which would
cause reducibility is therefore impossible.
The rank-three parameter over $F$ comes from a cuspidal
strong base change: the residual genericity criterion
after Thorne's Corollary 4.3 excludes noncuspidal discrete
blocks, and its global irreducibility excludes multiple
cuspidal summands in Theorem 4.2.
Cuspidal cyclic base change and Thorne's Theorem 4.4 now
give the stated actual definite-unitary representation on
$H/H^+$, with its exact split local components and weight.
All ramification remains relatively split, and the ordered
ordinary slopes and parallel weights restrict compatibly.

\emph{8. The previously certified rank bound on this actual tower.}
All quadratic factors are already in $H_{\rm ab}=F_1E$.
The first step in (G120.9) is cyclic degree $p$, and the
second is also cyclic degree $p$: $T_\theta$ ramifies only
at $q_0$, whereas $D_q/E$ ramifies only at $q\ne q_0$.
Both steps are unramified at $p$, while the cyclotomic
$\mathbb Z_p$-tower of the unramified base local fields is
totally ramified there.  Hence each finite step is disjoint
from the corresponding cyclotomic tower.
The residual coefficient splitting extension has degree
prime to $p$, giving the commuting product groups required
in AC-G100.  The local cubic inertia remains nontrivial,
so the permutation correction in AC-G74 is zero.
AC-G103's unconditional four-sector certificate gives
$r_{\alpha,H_{\rm ab}}\le a_{{\rm ab},0}\le3$ and
$[F_{\rm ab}:\mathbb Q]=4$.
Applying AC-G100 at the two actual cyclic-$p$ steps yields
$r_{\alpha,H}\le3p^2$, which is (G120.10).
For rational central and trivial coefficients, the real
host after adjoining their abelian splitting fields and
$\mu_p$ is a finite $p$-extension of an abelian field.
The Ferrero--Washington plus finite-$p$ ascent argument
already proved in AC-G76/AC-G98 consequently remains valid.
Thus subsequent AC-G92/AC-G93-fill and cyclotomic operations
retain the normalized strict-rank bound and their explicitly
proved cohomological consequences.
No arithmetic estimate for a new coefficient sector has
been manufactured from the automorphic construction.

For (G120.9a), the fixed field of the normal $p$-Sylow
subgroup of $\operatorname{Gal}(H/\mathbb Q)$ is
$H_{\rm ab}$: this subgroup is the product of
the rotations of $D_q/E$ and $\operatorname{Gal}(T_\theta/\mathbb Q)$.
Thus $H\cap\Omega\subset H_{\rm ab}$.
The three quadratic subextensions of $H_{\rm ab}/K$
are obtained by adjoining $\sqrt{-23},\sqrt5$, and
$\sqrt{-115}$.  Each ramifies at $23$ or $5$ over $K$,
so none lies in $\Omega$; this proves the equality.
Consequently the scalar auxiliary-place construction of
AC-G89 has the same actual finite-field compatibility on
$H$ as on the original quadratic host.  The native
$\mathrm{Lev}$, point and support applications still
require their literal records to be instantiated; they
are not assertions about an original-level $\mathrm{Dat}$.
\end{proof}

\paragraph{Source and scope record.}
For this extension the complete printed proof of Saha
Section 6.1.1/Theorem 6.2, pp.~901--902, and the complete
applied proof of Thorne Theorem 4.5, pp.~21--23, together
with Proposition 6.5 and its entire short proof, pp.~38--39,
were freshly read.  Thorne Corollary 4.3 and Theorem 4.4
are used at their exact stated scope.  Their deep transfer,
purity and local--global inputs remain published inputs.
The full-unit Jacquet and filtered-module calculation is
the independently checked AC-G79 proof, with its local
degree explicitly recomputed here.  Clozel--Thorne I,
Definition 2.4, Lemmas 2.5--2.6 and their complete printed
proofs, and Theorem 2.7 with its complete short proof,
were also freshly read through the cached primary PDF
retrieval; a direct PDF download returned 503, and no
local PDF byte identity is claimed.  The exact primary
excerpt and access scope are recorded alongside this file.
The actual universal-module, four-test, finite-character,
local-type and ray-cover arguments used above are the
already proved AC-G72, AC-G75, AC-G78--AC-G84, with the
changed arithmetic hypotheses verified in this proof.
No theorem saying that an arbitrary auxiliary tame-level
form satisfies the original $\mathrm{Dat.Kill}$ is used.
The two fields $T_\theta$ and $D_q$ have different roles:
the first removes the Eisenstein auxiliary type before
the unitary construction, and the second removes its
new level-raising type afterward.
